\documentclass{beamer}
\usepackage{amsfonts,amsmath,oldgerm}
\usetheme{sintef}
\usepackage{xeCJK}

% ---- 中文字体（Windows 自带，避免回退到 Fandol）----
\setCJKmainfont{Microsoft YaHei}
\setCJKsansfont{Microsoft YaHei}
\setCJKmonofont{SimHei}

\usepackage{booktabs}
\usepackage{array}
\usepackage{tikz}
\usetikzlibrary{positioning}

\newcommand{\testcolor}[1]{\colorbox{#1}{\textcolor{#1}{test}}~\texttt{#1}}

\usefonttheme[onlymath]{serif}

\titlebackground*{assets/background}

\newcommand{\hrefcol}[2]{\textcolor{cyan}{\href{#1}{#2}}}

\title{函数漏洞检测智能体工具}
\subtitle{课程项目开题报告}
\author{组长：李柏谚\quad 组员：杨聃、王义卓、孙伊凡、徐培栋}
\date{2026 秋季学期 · 高级软件工程}

\begin{document}
\maketitle

% ===================== 1. 项目意义 =====================
\footlinecolor{sintefgreen}
\begin{frame}{一、项目意义：我们为什么要做这件事}
\framesubtitle{漏洞太多、老工具不好用，而大模型还缺一点「安全常识」}
\small
\begin{itemize}
\item \textbf{漏洞多到查不过来}：2026 年前两个季度新增的公开漏洞就超过 \alert{3.5 万个}
      （每个漏洞有一个 CVE 编号，记录成因、影响和修补办法），人工审计根本跟不上。
\item \textbf{老工具不太好使}：传统静态分析工具靠写死的规则去匹配，
      规则没覆盖就\textbf{漏报}、规则太宽就\textbf{误报}，还说不出「为什么觉得有问题」。
\item \textbf{大模型有机会，但缺两样东西}：它\emph{读得懂代码}，
      但不知道具体漏洞长什么样（缺「安全常识」），也不知道谁调用了这个函数（缺「上下文」）。
\end{itemize}
\begin{colorblock}[black]{sinteflightgreen}{本项目的创新点}
\begin{itemize}
\item \textbf{不急着下结论}：像查案一样，先找证据、再说结论
\item \textbf{两路取证}：既翻「漏洞案例手册」，也去仓库查调用关系
\item \textbf{全程留痕}：查了什么、读了什么、为什么这么判，都在界面上
\end{itemize}
\end{colorblock}
\end{frame}

% ===================== 2. 项目内容 =====================
\begin{frame}{二、项目内容：我们要做一个什么东西}
\framesubtitle{输入一段代码，它告诉你「有没有漏洞、为什么、怎么查出来的」}
\small
\begin{block}{一句话说清楚}
给智能体一段 \textbf{C/C++ 函数代码} 和它所在\textbf{代码仓库的地址}，
它回答三件事：\alert{有没有漏洞} —— \alert{凭什么这么说} —— \alert{它是怎么查出来的}。
\end{block}
\begin{columns}[t]
\begin{column}{0.48\textwidth}
\begin{colorblock}[black]{sinteflightgreen}{难点一：大模型不懂安全知识}
\textbf{人话版}：它像个聪明但不熟悉业务的实习生，得先给它一本「漏洞案例手册」。\\
\textbf{做法}：把历史上已披露的真实漏洞案例整理成可检索的知识库，每次检测先去翻这本手册，找相似的旧案例作参考。
\end{colorblock}
\end{column}
\begin{column}{0.48\textwidth}
\begin{colorblock}[white]{sinteflilla}{难点二：光看一个函数判断不了}
\textbf{人话版}：只看一页纸，判断不了整件事——得知道它被谁调用、又调用了谁。\\
\textbf{做法}：给智能体装上「去仓库里查资料」的本事，让它自己决定什么时候需要查、查什么。
\end{colorblock}
\end{column}
\end{columns}
\vspace{0.1em}
{\footnotesize \textbf{我们主动把范围收窄}：只做「判断 + 解释」，不做自动修复、不扫全仓库。先把流程跑通，再谈性能。}
\end{frame}

% ===================== 3. 需求建模（业务 / 用户） =====================
\begin{frame}{三、需求建模（1）业务需求与用户需求}
\framesubtitle{谁会用、在什么场合用、用来解决什么问题}
\small
\begin{block}{业务需求}
在代码上线之前，快速把「\textbf{看着可能有漏洞}的函数」筛出来，
让安全人员集中精力看真正可疑的地方，而不是把代码从头看到尾。
\end{block}
\begin{columns}[t]
\begin{column}{0.5\textwidth}
\begin{colorblock}[white]{sintefpurple}{用户需求：谁会用它}
\begin{itemize}
\item \textbf{安全工程师}：日常代码审计，成批排查
\item \textbf{开发人员}：提交代码前自己先查一遍
\item \textbf{学生 / 研究者}：了解漏洞模式，做对比实验
\end{itemize}
\end{colorblock}
\end{column}
\begin{column}{0.5\textwidth}
\begin{block}{典型使用场景}
\begin{enumerate}
\item 把一段函数和仓库地址贴进界面
\item 界面实时显示它每一步在查什么
\item 拿到结论：有无漏洞 + 理由 + 分析记录
\item 需要时上传文件，一次跑几十上百条
\end{enumerate}
\end{block}
\end{column}
\end{columns}
\end{frame}

% ===================== 4. 需求建模（功能） =====================
\begin{frame}{三、需求建模（2）功能需求}
\framesubtitle{系统必须能做的事，一共六项}
\centering
\footnotesize
\begin{tabular}{@{}p{0.24\textwidth}p{0.68\textwidth}@{}}
\toprule
\textbf{功能} & \textbf{说明} \\
\midrule
F1 单样本检测 & 贴一段函数 + 仓库地址，给出「有漏洞 / 无漏洞」的结论 \\
F2 批量检测     & 上传 \texttt{.jsonl} 文件一次跑多条，汇总生成结果文件 \\
F3 知识检索     & 自动去「漏洞案例手册」里找相似案例，并展示找到了哪几条 \\
F4 上下文挖掘   & 自动去代码仓库查调用方、被调用方、全局变量、类型与宏 \\
F5 报告生成     & 输出「判定结果 + 判定理由 + 分析记录」三件套 \\
F6 过程可视化   & 界面实时展示它每一步在做什么，事后还能回看整个分析过程 \\
\bottomrule
\end{tabular}
\vspace{0.6em}

{\footnotesize 六项功能围绕一条主线：\emph{看得见的检测过程}——不仅要给结论，还要让人能复核这个结论。}
\end{frame}

% ===================== 5. 需求建模（非功能） =====================
\begin{frame}{三、需求建模（3）非功能需求}
\framesubtitle{做得好不好，看这五条}
\small
\begin{itemize}
\item \textbf{判断要准}：核心考核标准是「同一个函数的\emph{漏洞版}和\emph{修复版}能否被区分开」，
      区分得越多分越高——而不是简单地「见谁都判有漏洞」。
\item \textbf{过程要透明}：每一次知识检索、每一次仓库读取都要留痕，结论可复核、可追问。
\item \textbf{合规要严格}（硬性红线）：一次只给一个版本、不给代码改动记录、不许联网搜答案。
\item \textbf{成本要可控}：限制工具调用次数与文件数、缓存仓库快照，避免跑一次耗时过久或触发限流。
\item \textbf{上手要容易}：\alert{必须有图形界面}，不接受只能敲命令行的版本。
\end{itemize}
\vspace{0.2em}
\begin{block}{一句话概括我们的取向}
宁要\emph{慢一点但说得清}的判断，也不要\emph{快但查不到依据}的结论——可解释性是这个项目的立足点。
\end{block}
\end{frame}

% ===================== 6. 系统架构 =====================
\begin{frame}{四、系统架构初步考虑}
\framesubtitle{四层结构：界面 $\rightarrow$ 大脑 $\rightarrow$ 双手 $\rightarrow$ 资料}
\centering
\begin{tikzpicture}[
  node distance=0.3cm,
  box/.style={draw=black!30, rounded corners=2pt, align=center,
              font=\footnotesize, minimum width=0.9\textwidth,
              minimum height=0.82cm, inner sep=3pt}
]
\node[box, fill=sinteflightgreen] (ui)
  {表现层（界面）：输入函数与仓库地址、实时看分析过程、接收报告};
\node[box, fill=sinteflilla, text=white, below=of ui] (agent)
  {智能体层（大脑）：想一想 $\rightarrow$ 用个工具 $\rightarrow$ 看看结果 $\rightarrow$ 再想，证据够了才下结论};
\node[box, fill=sintefyellow, below=of agent] (tools)
  {工具层（双手）：\ ① 翻漏洞案例手册\ \ $|$\ \ ② 去代码仓库查资料（grep / 找调用关系）};
\node[box, fill=maincolor, text=white, below=of tools] (data)
  {数据层（资料）：历史漏洞案例库\ \ $|$\ \ 仓库快照缓存};
\end{tikzpicture}
\vspace{0.4em}

{\footnotesize 分层的好处：\textbf{界面、大脑、双手、资料各管一摊}，可以分开开发和替换——比如以后换一个更聪明的大模型，其余三层不用动。}
\end{frame}

% ===================== 7. 技术栈 =====================
\begin{frame}{五、开发技术栈与工具平台}
\framesubtitle{原则：选成熟、资料多、踩坑少的技术，把力气花在真正的难点上}
\centering
\footnotesize
\renewcommand{\arraystretch}{1.15}
\begin{tabular}{@{}p{0.17\textwidth}p{0.28\textwidth}p{0.47\textwidth}@{}}
\toprule
\textbf{模块} & \textbf{选型} & \textbf{为什么选它} \\
\midrule
开发语言 & Python 3.11 & 生态最全，前后端与算法一把梭 \\
智能体框架 & LangChain（Agent + 工具调用） & 成熟、示例最多，「大脑思考」的架子现成 \\
知识库 / 检索 & Chroma 向量库 + 中文嵌入模型 & 本地就能跑、零服务依赖，便于复现 \\
代码理解 & GitHub API + Tree-sitter + grep/find & 精确拉取指定版本代码，并画出调用关系 \\
图形界面 & Streamlit & 天然支持实时刷新与文件上传，开发量小 \\
底层大模型 & 兼容 OpenAI 协议的大模型 & 接口统一，随时可换模型做对比 \\
协作与版本 & Git + GitHub（Issue / 分支 / PR） & 过程留痕，团队协作数据可查 \\
开发工具 & VS Code / PyCharm & 免费、插件多，团队通用 \\
\bottomrule
\end{tabular}
\end{frame}

% ===================== 8. Git 协作 =====================
\begin{frame}{六、Git 协作开发：过程数据可查}
\framesubtitle{五个人的代码怎么合到一起，且每一步都留得下痕迹}
\footnotesize
\begin{columns}[t]
\begin{column}{0.52\textwidth}
\begin{block}{我们的工作流}
\begin{itemize}
\item \textbf{仓库}：GitHub 建团队仓库，代码、文档、脚本统一托管
\item \textbf{分支}：\texttt{main}（随时可演示）$\leftarrow$ \texttt{dev}（集成）$\leftarrow$ \texttt{feat/xxx}（各人开发）
\item \textbf{任务}：每项任务先开 Issue，认领后动手，看板上看得到进度
\item \textbf{合并}：一律走 Pull Request，至少一人 review 后才合并
\item \textbf{提交}：commit 写清「做了什么」，不写「update」
\end{itemize}
\end{block}
\end{column}
\begin{column}{0.46\textwidth}
\begin{colorblock}[white]{sintefpurple}{为什么强调「可查」}
\begin{itemize}
\item 提交记录、Issue、PR、评审意见全部公开
\item 谁做了什么、做到哪一步，一目了然
\item 既是协作依据，也是个人贡献的凭证
\item 出问题能回溯到具体哪次改动
\end{itemize}
\end{colorblock}
\end{column}
\end{columns}
\end{frame}

% ===================== 9. 进度安排 =====================
\begin{frame}{七、项目开发计划与进度安排}
\framesubtitle{七周里程碑，每周都有能拿得出手的东西}
\centering
\footnotesize
\begin{tabular}{@{}p{0.07\textwidth}p{0.25\textwidth}p{0.6\textwidth}@{}}
\toprule
\textbf{周次} & \textbf{阶段} & \textbf{本周产出（可验收）} \\
\midrule
W1 & 需求与数据摸底 & 理清数据字段；手工跑通几个样例，定下评分口径 \\
W2 & 搭建知识库     & 历史案例入库，检索功能可用并能自测效果 \\
W3 & 仓库查询工具   & 能拉取指定版本代码，能解析函数调用关系 \\
W4 & 打通智能体     & 「先检索 $\rightarrow$ 再判断 $\rightarrow$ 需要时查仓库」的完整流程跑通 \\
W5 & 界面与批量     & 图形界面可用（输入 / 过程 / 报告）+ 批量文件检测 \\
W6 & 联调调优与评测 & 300 条一起跑，按评价标准统计并迭代优化 \\
W7 & 报告与验收     & 技术报告、系统演示、结果文件提交 \\
\bottomrule
\end{tabular}
\vspace{0.5em}

{\footnotesize \textbf{节奏}：每周一次组会同步进度；从 W4 开始，保证「任何时候都能当场演示一个完整流程」。}
\end{frame}

% ===================== 10. 团队分工 =====================
\begin{frame}{八、项目团队组成与工作安排}
\framesubtitle{五人各管一摊，接口先定好再并行开发}
\small
\begin{itemize}
\item \textbf{李柏谚（组长）}：总体架构与智能体主流程；把控进度、协调接口
\item \textbf{杨聃}：历史漏洞案例的知识库与检索（把「手册」做出来并保证查得准）
\item \textbf{王义卓}：仓库查询工具（拉取代码、解析调用关系、封装检索命令）
\item \textbf{孙伊凡}：图形界面与批量流程（过程可视化、结果展示）
\item \textbf{徐培栋}：评测、集成与文档（跑批统计、回归测试、报告撰写）
\end{itemize}
\vspace{0.2em}
\begin{block}{协作方式}
数据接口先行——先把「一条案例长什么样、工具返回什么、报告有哪些字段」定死，
各人再并行开发，最后统一集成；全员参与联调与评测。
\end{block}
\end{frame}

% ===================== 11. 预期成果 =====================
\begin{frame}{九、预期成果与评价目标}
\framesubtitle{最后交付什么，以及我们打算拿到什么水平}
\footnotesize
\begin{columns}[t]
\begin{column}{0.48\textwidth}
\begin{block}{交付物}
\begin{itemize}
\item 一个能跑的完整系统：图形界面 + 知识库 + 仓库查询工具
\item 300 条样本的检测结果文件
\item 技术报告与现场演示
\end{itemize}
\end{block}
\begin{colorblock}[black]{sinteflightgreen}{我们的目标}
核心看「\textbf{同一函数的漏洞版与修复版能否被正确区分}」：
\textbf{争取 B 级及以上}（区分对 \textbf{25 组以上}），
同时把「\textbf{完全判反}」的情况压到区分对数的一半以内。
\end{colorblock}
\end{column}
\begin{column}{0.48\textwidth}
\begin{block}{可能踩的坑与预案}
\begin{itemize}
\item \textbf{把修复版也判成有漏洞}：提高门槛，必须能指出具体可疑点才判「有漏洞」
\item \textbf{把漏洞版判成没问题}：加强案例手册的提示作用，放宽初筛
\item \textbf{接口被限流}：仓库代码本地缓存，不重复拉取
\item \textbf{进度紧张}：先保主流程能演示，界面细节往后放
\end{itemize}
\end{block}
\end{column}
\end{columns}
\end{frame}

\end{document}
